\documentclass[12pt]{article}

\usepackage[a4paper,  top=1.3in, bottom=1.4in, left=1.4in, right=1.5in]{geometry}

\usepackage[utf8]{inputenc}
\usepackage[T1]{fontenc}
\usepackage{amsmath}
\usepackage{amsfonts}
\usepackage{amssymb}
\usepackage{tabularx}
\usepackage{array}
\usepackage{float}

\usepackage{graphicx, float}
\usepackage{adjustbox}
\graphicspath{{images/}}

\renewcommand{\figurename}{Slika}
\renewcommand{\tablename}{Tabela}


\renewcommand{\baselinestretch}{1.2} % Line spacing

% Parskip and parindent
\setlength{\parindent}{0pt} % Begin of paragraph indentation
\setlength{\parskip}{1em} % Paragraph spacing

%--- for references ---------
\usepackage[
    backend=biber,
    style=numeric,
    sorting=none
    ]{biblatex}

\addbibresource{bibliography.bib}
\AtBeginBibliography{\vspace*{10pt}}
%----------------------------

%----------- toc ------------------------
\usepackage{tocloft}
\usepackage{fancyhdr}
% \setlength{\cftbeforesecskip}{0pt} % Adjust the spacing before section entries
% \setlength{\cftbeforesubsecskip}{0pt} % Adjust the spacing before subsection entries
%----------------------------------------

\begin{document}
\pagenumbering{gobble}  % Suppress page numbers initially

% Title Page
\newgeometry{top=1in, bottom=1in, left=1in, right=1in} % New margins for title page
\begin{titlepage}
    \begin{center}
        
        % add your university logo here
        % negative value moves the logo up
        \vspace*{-1in}
        \includegraphics[width=0.4\textwidth]{raf_logo.png}

        % set font size to 14pt
        \vspace{1in}
        \Large
        \textbf{Tema}
        
        % set horizontal margin for the title to 1.5in and center it
        \vspace{1in}
        \Huge
        \textbf{Skaliranje sistema za rezervaciju: Analiza različitih arhitekturalnih pristupa i strategija za skaliranje servisa}
        
        \vspace{1in}


            \fontsize{17pt}{17pt}\selectfont
            \textbf{Kurs: Praktikum in računarstva u oblaku} \\
            \vspace*{1.5in}
            
            \begin{center}
            \normalsize
            \begin{tabular}{p{0.75\textwidth} p{0.5\textwidth}}
                \fontsize{14pt}{18pt}\selectfont   
                \textbf{Mentor:} & 
            
                \fontsize{14pt}{18pt}\selectfont
                \textbf{Student:} \\
                prof. Mirjana Radivojević & Vanja Kovinić \\
            \end{tabular}
            \end{center}

            \vspace*{\fill}

            \normalsize
            Beograd, 2025.


            
        \end{center}
    \end{titlepage}
    \restoregeometry % Restore original margins

    \newpage
    \newgeometry{top=1.3in, bottom=2.2in, left=1.4in, right=1.4in} % New margins for title page
    \renewcommand{\contentsname}{Sadržaj}
    \tableofcontents{}

    \newpage
    \clearpage             % Move to a new page
    \pagenumbering{arabic} % Start using Arabic numerals
    \setcounter{page}{1}   % Set the counter to 1

    \newpage
    \section{Uvod}
    \subsection{Problem i motivacija}

    Upravljanje resursima u modernim poslovnim okruženjima predstavlja ključni izazov za 
    efikasno funkcionisanje organizacija. Rezervacija sala za sastanke, kao jedan od najčešćih 
    operativnih procesa, često je opterećena manuelnim postupcima koji dovode do konflikata 
    u rasporedu, duplih rezervacija i neoptimalnog korišćenja prostora.

    U konkretnom slučaju koji je analiziran u ovom radu, vlasnik poslovnog centra koji 
    iznajmljuje sale za sastanke suočavao se sa potpunim odsustvom digitalnog sistema 
    za rezervacije. Ceo proces rezervacije bio je centralizovan kroz jednu sekretaricu koja 
    je putem email komunikacije primala zahteve, proveravala dostupnost i potvrđivala 
    rezervacije. Ovakav pristup stvorio je više kritičnih problema:

    \begin{itemize}
        \item \textbf{Nedovoljna preglednost}: Klijenti nisu imali mogućnost da vide raspoloživost sala u realnom vremenu, što je otežavalo planiranje.
        \item \textbf{Ograničena fleksibilnost}: Promene u rasporedu bile su teške za implementaciju, što je dovodilo do dodatnih konflikata.
        \item \textbf{Neproduktivno trošenje vremena}: Značajan deo radnog vremena sekretarice bio je posvećen repetitivnim zadacima umesto obavljanju kompleksnih i strateških zadataka.
        \item \textbf{Neprofesionalan imidž}: Manuelni proces ostavljao je utisak zastarele i neorganizovane kompanije kod klijenata.
        \item \textbf{Nedostatak analitike}: Bez digitalnog sistema, bilo je nemoguće pratiti trendove korišćenja sala, što je otežavalo donošenje poslovnih odluka.
        \item \textbf{Skalabilnost problema}: Sa rastom broja klijenata, sistem je postajao sve manje održiv.
    \end{itemize}

    \subsection{Tehnologije i pristup}

    Implementacija sistema zasniva se na \textbf{Java} \cite{java_language} programskom jeziku 
    i \textbf{Spring Boot} \cite{spring_boot} 
    okruženju, koje čini osnovni sloj aplikacione logike. Uz savremene pristupe 
    kontejnerizacije i principe razvoja \textit{cloud-native} aplikacija, postignuta je 
    visoka prenosivost i konzistentnost sistema u različitim okruženjima. \textbf{Docker} \cite{docker} je 
    korišćen za pakovanje i izolaciju aplikacije, dok je \textbf{GitHub Actions} \cite{github} iskorišćen 
    za potpunu automatizaciju procesa kontinuirane integracije i isporuke (\textbf{\textit{CI/CD}} \footnote{\textit{CI/CD - Continuous Integration/Continuous Deployment}}).

    Aplikacija je postavljena na \textbf{Netcup cloud} \cite{netcup} infrastrukturu, čime je obezbeđeno ekonomično i 
    skalabilno rešenje za produkcijski hosting.

    Ključne tehnologije korišćene u projektu su:
    \begin{itemize}
    \item \textbf{Razvojna platforma:} \textbf{Java} i \textbf{Spring Boot}, sa implementacijom \textit{RESTful API}-ja\footnote{\textit{RESTful API (Representational State Transfer)} je pristup pravljenja web servisa koji koriste \textit{HTTP} protokol za komunikaciju i manipulisanje resursima na serveru.}
 za komunikaciju između klijenta i servera
    \item \textbf{Kontejnerizacija:} \textbf{Docker}, za pakovanje aplikacije i izolaciju od okruženja  
    \item \textbf{Orkestracija:} \textbf{Docker Compose}, za upravljanje više kontejnera u okviru jednog sistema  
    \item \textbf{Automatizacija (CI/CD):} \textbf{GitHub Actions}, za automatsko testiranje i postavljanje aplikacije na server  
    \item \textbf{Baza podataka:} Relaciona baza podataka (\textbf{MySQL} \cite{mysql}) pokrenuta u posebnom kontejneru  
    \item \textbf{Cloud infrastruktura:} \textbf{Netcup VPS}, korišćen za postavljanje i pokretanje sistema u produkciji
    \end{itemize}
    
    Poseban akcenat stavljen je na automatizaciju procesa postavljanja aplikacije, 
    što omogućava brzo isporučivanje novih funkcionalnosti kroz jednostavno slanje izmena u 
    repozitorijum (\textit{git push}). Ovakav pristup značajno smanjuje vreme potrebno za 
    implementaciju i testiranje izmena.
    Izvorni kod i konfiguracija dostupni su na \textbf{GitHub} platformi \footnote{\url{https://github.com/Kovelja009/conf_room}}.

    \section{Arhitekturni obrasci: Monolit vs Mikroservisi}
    \subsection{Teoretske osnove}

    \textbf{Monolitna arhitektura} predstavlja tradicionalni pristup razvoju softverskih sistema
    , pri kojem se celokupna funkcionalnost implementira u okviru jedne celovite i 
    zajednički postavljene aplikacije. U ovom modelu, sve komponente sistema - upravljanje 
    korisnicima, poslovna logika, pristup bazi podataka i prikaz podataka 
    (interfejs) - integrisane su unutar jednog izvršnog procesa. Komunikacija između delova 
    sistema odvija se putem direktnih poziva funkcija unutar iste aplikacije, što može 
    doprineti boljim performansama, ali istovremeno dovodi do 
    jake povezanosti (\textit{tight coupling}) između modula i smanjenja fleksibilnosti u 
    razvoju i održavanju.

    S druge strane, \textbf{mikroservisna arhitektura} predstavlja savremeniji pristup 
    koji uvodi koncept distribuiranih sistema. Aplikacija se u ovom modelu sastoji od niza 
    nezavisnih servisa, od kojih je svaki odgovoran za tačno određenu poslovnu funkcionalnost. 
    Svaki mikroservis poseduje sopstvenu bazu podataka i može se razvijati, testirati i 
    postavljati potpuno nezavisno od ostalih servisa. Ovakav pristup omogućava veću fleksibilnost 
    i skalabilnost, ali istovremeno uvodi dodatnu složenost, naročito u oblastima kao što su 
    otkrivanje servisa (\textit{service discovery}), međuservisna komunikacija i upravljanje 
    transakcijama koje obuhvataju više servisa.

    \newpage
    U kontekstu razvoja za \textit{cloud} okruženje, oba arhitekturna pristupa mogu se 
    uspešno implementirati uz korišćenje kontejnerskih tehnologija. 
    \textbf{Docker} omogućava konzistentno postavljanje aplikacija bez obzira na konkretno 
    okruženje u kojem se izvršavaju, dok platforme za orkestraciju poput 
    \textbf{Kubernetesa} \cite{kubernetes} nude napredne mogućnosti za upravljanje servisima, automatsko 
    skaliranje i raspodelu opterećenja. Konačan izbor arhitekture zavisi od specifičnih zahteva 
    aplikacije, veličine i organizacije razvojnog tima, kao i infrastrukturnih i operativnih 
    ograničenja.
    
    \newpage
    \subsection{Komparativna analiza}

    Poređenje monolitne i mikroservisne arhitekture prikazano je u Tabeli \ref{tabela:monolit_vs_mikroservisi}. 
    Ova tabela sumira ključne aspekte oba pristupa, uključujući skalabilnost, 
    izolaciju grešaka, kompleksnost razvoja, performanse, složenost \textit{CI/CD} procesa i 
    tehnologije koje se koriste. 

    \hspace{0.4cm}
    \begin{table}[H]    
    \centering
    \small
    \renewcommand{\arraystretch}{1.4}
    \begin{tabularx}{\textwidth}{
        |>{\raggedright\arraybackslash}>{\hsize=0.8\hsize}X
        |>{\raggedright\arraybackslash}>{\hsize=0.8\hsize}X
        |>{\raggedright\arraybackslash}>{\hsize=0.8\hsize}X
        |>{\raggedright\arraybackslash}>{\hsize=1.6\hsize}X|}
            \hline
            \textbf{Aspekt} & \textbf{Monolit} & \textbf{Mikroservisi} & \textbf{Kontekst primene} \\
            \hline
            Skalabilnost & Vertikalna & Horizontalna & Monolit za mali broj korisnika, mikroservisi za veći broj i opterećenje \\
            \hline
            Izolacija grešaka & Niska & Visoka & Manji sistemi mogu tolerisati jednu tačku otkaza \\
            \hline
            Kompleksnost razvoja & Niska & Visoka & Monolit za male timove, mikroservisi za kompleksne sisteme \\
            \hline
            Performanse & Brže (lokalni pozivi) & Sporije (mreža) & Direktan metod vs \textit{HTTP/REST} \\
            \hline
            \textit{Deployment} & Jednostavan & Složen & \textit{CI/CD} zahteva više skripti za mikroservise \\
            \hline
            Komunikacija & Lokalni pozivi & Međuservisna komunikacija & Direktan pristup vs API pozivi \\
            \hline
            \textit{CI/CD} složenost & Jedan tok & Više tokova & Veći broj servisa, nezavisna postavljanja \\
            \hline
            Tehnologije & Jedinstven stek & Potencijalno više tehnologija & Manji timovi preferiraju doslednost \\
            \hline
            Monitoring & Centralizovan & Distribuiran & Lakše praćenje u monolitu \\
            \hline
            \end{tabularx}
            \caption{Uporedna analiza monolitne i mikroservisne arhitekture}
        \label{tabela:monolit_vs_mikroservisi}


    \end{table}
    \newpage

    \subsection{Faktori za donošenje odluke}

    Izbor arhitekture samog servisa zavisi od više povezanih faktora koje je potrebno pažljivo 
    analizirati u konkretnom poslovnom i tehničkom kontekstu:

    \subsubsection*{Struktura tima i organizacioni faktori}

    \textit{Conway-ev zakon} sugeriše da arhitektura sistema odražava strukturu komunikacije 
    unutar organizacije koja ga razvija. U slučaju malih timova (npr. dva programera), 
    što je karakteristično za analizirani projekat, \textbf{monolitna arhitektura} se često 
    pokazuje kao optimalan izbor jer omogućava:
    \begin{itemize}
        \item jedinstveno razvojno okruženje bez dodatnog koordinacionog opterećenja između \textit{frontend} i \textit{backend} programera,
        \item pojednostavljeno orkestriranje koda - direktan pristup zajedničkom kodu,
        \item brže donošenje odluka bez zavisnosti između timova i dileme oko vlasništva nad servisima.
    \end{itemize}

    \subsubsection*{Količina saobraćaja i zahtevi u pogledu performansi}

    Očekivano opterećenje predstavlja ključni faktor u izboru arhitekture:
    \begin{itemize}
        \item \textbf{Nizak nivo saobraćaja} (desetine istovremenih korisnika): performanse monolita su 
        u potpunosti dovoljne, dok bi mrežna latencija karakteristična za mikroservise predstavljala 
        nepotreban trošak.
        \item \textbf{Srednji nivo saobraćaja} (stotine korisnika): prelazna zona u kojoj su oba 
        pristupa moguća i zavise od dodatnih faktora.
        \item \textbf{Visok nivo saobraćaja} (hiljade i više korisnika): mikroservisna arhitektura 
        omogućava horizontalno skaliranje i bolje upravljanje opterećenjem, što je ključno za 
        očuvanje performansi i dostupnosti sistema.
    \end{itemize}

    \subsubsection*{Brzina razvoja i učestalost \textit{deployment}-a aplikacije}

    Učestalost \textit{deployment}-a značajno utiče na odabir arhitekture:
    \begin{itemize}
        \item \textbf{Čest \textit{deployment}} (npr. više puta dnevno tokom razvoja): automatizacija 
        \textit{CI/CD} procesa postaje od suštinskog značaja, ali je jednostavnija za 
        implementaciju u okviru monolita.
        \item \textbf{Koordinacija više servisa}: mikroservisi zahtevaju naprednu orkestraciju kako 
        bi se obezbedila sinhronizovana postavljanja.
    \end{itemize}

    \subsubsection*{Tehnološka ograničenja i operativna zrelost}

    DevOps kapaciteti organizacije predstavljaju značajan ograničavajući faktor:
    \begin{itemize}
        \item \textbf{Ograničeno operativno iskustvo}: postavljanje i nadgledanje monolita je 
        jednostavnije i zahteva manje alata i znanja.
        \item \textbf{Napredni operativni kapaciteti}: mikroservisi zahtevaju složenije mehanizme 
        za otkrivanje servisa (\textit{service discovery}), balansiranje 
        opterećenja (\textit{load balancing}) i distribuirani nadzor sistema.
        \item \textbf{Kompleksnost komunikacije između servisa}: direktni pozivi metoda naspram 
        \textit{HTTP/REST API} poziva za razmenu podataka.
    \end{itemize}

    Konačna odluka treba da uspostavi ravnotežu između trenutnih potreba za brzinom razvoja i 
    dugoročnih zahteva za skalabilnošću, uzimajući u obzir sposobnosti postojećeg tima i 
    predviđeni rast sistema.



\newpage
\printbibliography[title={Literatura}]
\end{document}